\documentclass[12pt]{article}
\usepackage[utf8]{inputenc}
\usepackage[T1]{fontenc}
\usepackage{hyperref}
\title{Kiorr Language Resources for K-12 Curriculum Development}
\author{Compiled by Claude Sonnet 5.5 (claude-sonnet-5-5) for Andr\'as Kornai}
\date{}
\begin{document}
\maketitle

\section{Introduction}
Kiorr (ISO 639-3 xko; Glottocode kior1239) is an Austroasiatic language of the Palaungic branch, spoken in Luang Namtha Province, Laos. Reported speaker figures differ: about 2,400 (a 1985 figure cited by Wikipedia) and roughly 4,700 to 5,200 (Joshua Project, Global Recordings Network). Other names include Kha Kior, Saamtaav and Samtao. Classification of the variety ``Con'' (cno) is disputed. The language appears to have no established written form. Searches found no school materials in or about Kiorr. Note that the survey's identification line gives the family as ``Bookkeeping'', which looks like a data error; the sources found place Kiorr in Austroasiatic.

\section{Primary Educational Resources}
No primers, textbooks or school programmes for Kiorr were found in the searches. Joshua Project describes the language as having no written form and no Scriptures or ministry materials. A K-12 curriculum would have to be built from scratch, starting with community-led orthography development.

\section{Linguistic and Reference Materials}
\begin{itemize}
\item Wikipedia, ``Kiorr language'': overview of classification, names and speaker estimates. \url{https://en.wikipedia.org/wiki/Kiorr_language}
\end{itemize}
No dictionary or grammar was found in the search results.

\section{Religious and Community Resources}
\begin{itemize}
\item Joshua Project, ``Kiorr in Laos'': people-group profile, population estimate, and the note that no Scriptures are available. \url{https://joshuaproject.net/people_groups/12706}
\item Joshua Project, language page for xko. \url{https://joshuaproject.net/languages/xko}
\end{itemize}

\section{Digital and Audio Resources}
\begin{itemize}
\item Global Recordings Network, language page for xko (listed in search results; the content was not inspected, so availability of recordings is unconfirmed). \url{https://globalrecordings.net/en/language/xko}
\end{itemize}

\section{Curriculum Development Context}
Very little documentation exists. Any curriculum would need community consultation, documentation fieldwork, an agreed orthography, and coordination with Lao national-language schooling. Because ``Kiorr'', ``Con'' and ``Samtao/Saamtaav'' may overlap, sources should be checked carefully so that materials from related varieties are not attributed to Kiorr.

\end{document}